% No \chapter here (the preface opens straight into the epigraph), so set the running
% head and TOC entry by hand. Without this, \leftmark keeps the stale "Table of Contents"
% mark and it leaks onto the preface pages.
\addcontentsline{toc}{chapter}{Preface}
\markboth{Preface}{}

% TODO: float me right
\vspace*{\fill}
    \noindent
    {\itshape\footnotesize 
       ``Take care what you're doing, Afranius. After all, they're temple seals!'' \\
       ``The procurator need not trouble himself about this question,'' replied Afranius, closing the packet. \\ ``You mean you have all their seals?'' asked Pilate, laughing. \\ 
       ``How could it be otherwise,'' replied Afranius dryly, with no trace of laughter. 
    }
    \vspace{2em}

    \noindent\hfill\hspace*{-1.5em}
    {{\large --- Mikhail Bulgakov} \\}
    \vspace{1em}
    \noindent\hfill\hspace*{-1.5em}
    {\textit{The Master and Margarita}\footnote{trans. Diana Burgin \& Katherine Tiernan O'Connor, pp. 275-276}}
\vspace*{\fill}

Borges once remarked that ``the generations of men, throughout recorded time, have always told and retold two stories---that of a lost ship which searches the Mediterranean seas for a dearly loved island, and that of a god who is crucified on Golgotha.''\footnote{
``The Gospel According to Mark'' from \emph{Dr. Brodie's Report}, trans. Norman Thomas di Giovanni, p. 9.}With regard to the latter, he was perhaps not entirely correct. The artistic and philosophical reinterpretation of the events surrounding the execution of Jesus Christ only became an established literary tradition in the 19\textsuperscript{th} century, when the Church lost a significant portion of its function of ideological supervision. Within a narrative, this theme typically takes the form of a subplot or a ``novel within a novel'' (as in the works of Bulgakov and Aitmatov); more rarely, it may take the form of a standalone work (such as those of France and Andreyev).\footnote{The works Yeskov is referencing are \emph{The Master and Margarita} by Mikhail Bulgakov, \emph{The Place of the
Skull} by Chingiz Aitmatov, ``The Procurator of Judea'' by Anatole France, and ``On the Day of the
Crucifixion'' by Leonid Andreyev. (Z.B.)} The texts that have been created within this tradition vary quite widely, both in their artistic level (from Bulgakov's immortal \emph{Master} to Ilya Varshavsky's comic sci-fi story ``Hysteresis Loop'') and in the degree to which they conform to Holy Scripture and historical reality (from the pinpoint precision of Yury Dombrovsky to the deliberate inaccuracy of the Strugatsky brothers).\footnote{These works are, respectively, \emph{The Faculty of Useless Knowledge and Burdened with Evil.} (Z.B.)} With regard to the latter aspect, it is worth comparing two renowned cinematic masterpieces: \emph{The Gospel According to St. Matthew} and \emph{The Last Temptation of Christ}. It goes without saying that all these authors' versions differ quite radically, to the point where the Gospel figures depicted in them are merely ``homonyms'' of each other---compare France's Pilate with Bulgakov's, Andreyev's Judas with Dombrovsky's, or Pasolini's Jesus with Scorcese's.

Nevertheless, there is still one fundamental constraint that all works produced within this tradition are subject to: supernatural forces may not directly interfere with the course of events other than to send ``a strange cloud drifting into Yershalaim.''\footnote{\emph{The Master and Margarita}, p. 255. ``Yershalaim'' is the Russian transliteration of the Hebrew word ``Yerushalayim,'' meaning ``Jerusalem.'' (Z.B.)} This is precisely why an event so crucial to the Christian worldview as the Bodily Resurrection is always omitted from these narratives---despite the fact that many of the authors who have turned to this theme were undoubtedly people of faith. And since I belong to a generation whose perceptions were shaped immeasurably more by Bulgakov's Yeshua\footnote{Similarly to his use of the word ``Yershalaim'' (see note 6), Bulgakov refers to Jesus in \emph{The Master and Margarita} by his Aramaic name, Yeshua Ha-Notsri. (Z.B.)} than by his official prototype, the problem of the Resurrection was one that did not interest me in the slightest up until very recently.